% !TEX root = ../../ctfp-print.tex

\lettrine[lhang=0.17]{范}{畴}（category）是一个极其简单的概念。
一个范畴由\newterm{对象}（object）和连接它们的\newterm{箭头}（arrow）组成。
这就是为什么范畴可以用图形直观表示的原因。
对象可以画成圆圈或点，箭头...就是箭头。（为了多样性，我会偶尔把对象画成小猪，箭头画成烟花。）
但范畴的本质在于\emph{组合性}（composition）。
或者换种说法：组合性的本质就是范畴。
箭头是可以组合的，因此如果你有一个从对象$A$到对象$B$的箭头，
另一个从对象$B$到对象$C$的箭头，
那么必然存在一个从$A$到$C$的箭头——它们的组合。

\begin{figure}
  \centering
  \includegraphics[width=0.8\textwidth]{images/img_1330.jpg}
  \caption{在范畴中，若存在从$A$到$B$的箭头和从$B$到$C$的箭头，
    则必须存在直接从$A$到$C$的组合箭头。
    这个图不是完整的范畴，因为它缺少恒等态射（见后文）。}
\end{figure}

\section{作为函数的箭头}

这已经太抽象了吗？不要绝望。
让我们具体化讨论。
把箭头（也称为\newterm{态射}（morphism））想象成函数。
你有一个函数$f$，它接受类型为$A$的参数并返回$B$；
还有一个函数$g$，它接受$B$并返回$C$。
你可以通过将$f$的结果传递给$g$来组合它们。
这样你就定义了一个新的函数：接受$A$并返回$C$。

在数学中，这种组合用函数间的小圆圈表示：$g \circ f$。
注意这里的右到左组合顺序。
对某些人来说这可能令人困惑。
你可能熟悉Unix中的管道符号：
\begin{snip}{text}
lsof | grep Chrome
\end{snip}
或者F\#中的右移运算符\code{>>}，它们都是从左到右的顺序。
但在数学和Haskell中，函数组合是右到左的。
如果你把$g \circ f$理解为「g \emph{在f之后}」会更容易理解。

让我们用C代码把这个过程变得更明确。
假设我们有一个函数\code{f}，它接受类型为\code{A}的参数并返回\code{B}：
\begin{snip}{text}
B f(A a);
\end{snip}
还有另一个函数：
\begin{snip}{text}
C g(B b);
\end{snip}
它们的组合是：
\begin{snip}{text}
C g_after_f(A a)
{
    return g(f(a));
}
\end{snip}
这里再次看到右到左的组合顺序：\code{g(f(a))}；这次是在C语言中。

我很想告诉你C++标准库中有模板可以接收两个函数并返回它们的组合，
但遗憾的是并没有这样的模板。
那我们换个话题看看Haskell。
这是从A到B的函数声明：
\src{snippet01}
同理：
\src{snippet02}
它们的组合是：
\src{snippet03}
当你见识过Haskell的简洁表达后，
就会觉得C++难以表达基本的函数式概念略显尴尬。
事实上，Haskell支持Unicode字符，
所以你可以这样写组合：
% don't 'mathify' this block
\begin{snip}{text}
g ◦ f
\end{snip}

甚至可以使用Unicode双冒号和箭头：
% don't 'mathify' this block
\begin{snipv}
f \ensuremath{\Colon} A → B
\end{snipv}
这是第一个Haskell知识点：双冒号表示「具有...类型」
函数类型通过在两个类型之间插入箭头创建。
通过插入句点（或Unicode圆圈）可以在两个函数间进行组合。

\section{组合的性质}

任何范畴中的组合都必须满足两个极其重要的性质：

\begin{enumerate}
  \item
        组合满足结合律。如果你有三个可以组合的态射 $f$、$g$ 和 $h$
        （即它们的对象首尾相接），组合时无需添加括号。用数学记号表示为：
        $$h \circ (g \circ f) = (h \circ g) \circ f = h \circ g \circ f$$
        在（伪）Haskell中表示为：

        \src{snippet04}[b]
        （我说"伪"是因为函数之间没有定义相等性）

        对于函数而言结合律显而易见，但在其他范畴中可能并不那么直观。

  \item
        每个对象 $A$ 都存在一个作为组合单位元的箭头。这个箭头从对象指向自身。
        所谓组合单位元是指：当与任意始于 $A$ 或终于 $A$ 的箭头组合时，
        结果都是原来的箭头。对象 A 的单位箭头 $\idarrow[A]$ 被称为 $A$ 上的
        \newterm{恒等箭头(identity)}。用数学记号表示，若 $f$ 是从 $A$ 到 $B$ 的态射，则
        $$f \circ \idarrow[A] = f$$
        且
        $$\idarrow[B] \circ f = f$$
\end{enumerate}
对于函数而言，恒等箭头通过返回其参数本身的恒等函数实现。该函数对所有类型都有相同的实现，
这意味着它是全称多态的。在C++中我们可以这样定义模板：

\begin{snip}{cpp}
template<class T> T id(T x) { return x; }
\end{snip}
当然，C++中并没有这么简单，因为你不仅要考虑传递什么，还要考虑传递方式
（按值传递、按引用传递、按常量引用传递、按移动传递等等）。

在Haskell中，恒等函数是标准库（Prelude）的一部分。以下是它的声明和定义：

\src{snippet05}
如你所见，Haskell中的多态函数非常简洁。在声明中只需将类型替换为类型变量即可。
这里有个约定：具体类型的名称总是以大写字母开头，类型变量名称以小写字母开头。
因此这里的 \code{a} 代表所有类型。

Haskell的函数定义由函数名后跟形式参数组成——此处只有一个参数 \code{x}。
函数体位于等号之后。这种简洁性常常让新手感到惊讶，但你会很快发现它非常合理。
函数定义和调用是函数式编程的核心要素，因此语法被简化到极致。不仅参数列表不需要括号，
参数之间也没有逗号（当你定义多个参数的函数时会看到这一点）。

函数体始终是一个表达式——函数中没有语句。函数的结果就是这个表达式——在此处就是 \code{x}。

这结束了我们的第二课Haskell教学。

恒等条件可以用（伪）Haskell写成：

\src{snippet06}
你可能会问自己：为什么有人要费心去定义什么都不做的恒等函数？同样的问题也可以问：
为什么我们要有零这个数字？零是"无"的符号。古罗马人使用的数字符号中没有零，但他们依然
建造了精良的道路和引水渠，其中许多至今仍在使用。

像零或 $\id$ 这样的中性元素在处理符号变量时极其有用。这就是为什么罗马人在代数方面表现欠佳，
而熟悉零概念的阿拉伯人和波斯人则擅长代数。因此恒等函数作为高阶函数的参数或返回值时特别方便。
高阶函数使得函数的符号操作成为可能，它们构成了函数的代数体系。

总结：范畴由对象和箭头（态射）构成。箭头可以组合，且组合满足结合律。每个对象都有一个恒等箭头作为组合的单位元。

\section{组合是编程的本质}

函数式程序员解决问题的方式颇为独特。他们通常会从一些颇具禅意的问题入手：
例如，在设计交互式程序时，他们会问：什么是交互？在实现康威生命游戏时，
他们可能会思考生命的真谛。本着这种精神，我也想提出一个问题：何谓编程？

从最基础的层面看，编程就是告诉计算机该做什么。「把内存地址x的内容取出，
加到寄存器EAX中去。」即使我们用汇编语言编程，给计算机的指令本质上都是
某种更有意义事物的表达。我们是在解决非平凡问题（如果是平凡问题，
我们就不需要计算机帮忙了）。那么我们如何解决问题呢？我们将大问题分解成小问题。
如果小问题仍然太大，就继续分解下去，直到最终写出能解决每个小问题的代码。
然后迎来编程的本质时刻：我们通过组合这些代码片段来构建更大问题的解决方案。
如果没有重组能力，分解也就失去了意义。

这种分层分解与重组的过程并非计算机强加给我们，而是反映了人类心智的局限性。
我们的大脑一次只能处理少量概念。心理学领域引用最广泛的论文之一——
\urlref{http://en.wikipedia.org/wiki/The_Magical_Number_Seven,_Plus_or_Minus_Two}{《神奇的数字7±2》}
提出，人类短期记忆只能容纳大约 $7 \pm 2$ 个「信息组块」。虽然我们对人脑短期记忆的理解在不断演变，
但确定无疑的是其容量有限。关键在于我们无法应对对象的混沌状态或代码的意大利面式纠缠。
我们需要结构不是因为结构优美的程序赏心悦目，而是因为缺乏结构会导致大脑无法高效处理。
我们常说某段代码优雅或优美，其实真正含义是它恰好适合人类有限的认知能力。
优雅的代码能创造出大小和数量都恰到好处的信息块，便于心智消化吸收。

那么程序组合的理想单元应该是什么样的呢？它们的表面积增长速度必须慢于体积。
（我喜欢这个类比，因为几何体的表面积随尺寸平方增长——比立方增长的体积要慢）
表面积是指组合这些单元所需的信息量，体积则是实现这些单元所需的信息量。
核心思想是：一旦某个单元完成实现，我们就可以忘记其实现细节，专注于它与其他单元的交互方式。
在面向对象编程中，表面积体现为对象的类声明或抽象接口；在函数式编程中，则体现为函数声明。
（这里做了简化，但基本原理如此）

范畴论在这方面走向极端：它明确反对我们窥视对象内部。范畴论中的对象是抽象模糊的实体，
你所能了解的全部信息就是它与其他对象的关系——即通过箭头建立的连接方式。
这类似于互联网搜索引擎通过分析网页的出入链接进行排名（当然作弊情况除外）。
在理想化的面向对象编程中，对象仅通过其抽象接口可见（纯表面积，无体积），
方法在此扮演箭头的角色。一旦你需要深入对象的实现细节才能理解如何组合使用它，
你就丧失了编程范式的固有优势。

\section{挑战}

\begin{enumerate}
  \tightlist
  \item
        在你最喜欢的编程语言（如果首选语言恰好是Haskell，则改用第二喜欢的语言）中尽可能地实现恒等函数(identity function)。
  \item
        在你最喜欢的编程语言中实现组合函数(composition function)。该函数接受两个函数作为参数，并返回它们的组合函数。
  \item
        编写一个程序，用于测试你的组合函数是否保持恒等性(respects identity)。
  \item
        从范畴论的角度来看，万维网(world-wide web)能否被视为一个范畴？超链接是否构成态射(morphisms)？
  \item
        能否将Facebook建模为范畴：以用户作为对象(objects)，好友关系作为态射(morphisms)？
  \item
        有向图在什么条件下能构成一个范畴？
\end{enumerate}